\documentclass[11pt,a4paper]{article}
\usepackage[utf8]{inputenc}
\usepackage[T1]{fontenc}
\usepackage[catalan]{babel}
\usepackage{lmodern}
\usepackage[sfdefault,lining]{FiraSans}
\usepackage[a4paper,top=2.3cm,bottom=2.2cm,left=1.8cm,right=1.8cm,headheight=1.3cm,headsep=0.4cm,footskip=1cm]{geometry}
\usepackage{amsmath,amssymb}
\usepackage{icomma}
\usepackage{graphicx}
\usepackage[table]{xcolor}
\usepackage{tikz}
\usepackage{fontawesome5}
\usepackage{ccicons}
\usepackage{fancyhdr}
\usepackage{lastpage}
\usepackage{tasks}
\usepackage{needspace}
\usepackage[most]{tcolorbox}
\usepackage{array}
\usepackage[hidelinks,pdftitle={Solucions de la fitxa de la unitat 10: Proporcionalitat. Teorema de Tales. 3r d'ESO},pdfsubject={Matemàtiques 3r ESO},pdfkeywords={raó, proporció, segments proporcionals, teorema de Tales, posició de Tales, semblança, raó de semblança, criteris de semblança, escala}]{hyperref}

\definecolor{navy}{HTML}{1B2748}
\definecolor{accent}{HTML}{E8603A}
\definecolor{teal}{HTML}{1F9E90}
\definecolor{soft}{HTML}{EEF1F8}
\definecolor{muted}{HTML}{5D6685}

\setlength{\parindent}{0pt}
\setlength{\parskip}{3pt}
\settasks{label=\alph*), label-format=\bfseries\color{navy}, label-width=1.3em, item-indent=1.8em, column-sep=0.6em, after-item-skip=5pt, before-skip=3pt, after-skip=4pt}

% ---------- capçalera i peu ----------
\newcommand{\llicencia}{\href{https://creativecommons.org/licenses/by-nc-sa/4.0/deed.ca}{\ccbyncsa}\ \ Llicència Creative Commons Reconeixement-NoComercial-CompartirIgual 4.0 Internacional (CC BY-NC-SA 4.0)}
\pagestyle{fancy}
\fancyhf{}
\fancyhead[L]{\raisebox{-0.35\height}{\includegraphics[height=0.95cm]{logo-gva.png}}}
\fancyhead[C]{\small\color{muted}Solucions de la fitxa · Unitat 10: Proporcionalitat. Teorema de Tales · 3r d'ESO}
\fancyhead[R]{\raisebox{-0.35\height}{\includegraphics[height=0.95cm]{logo-vermellar.png}}}
\fancyfoot[L]{\scriptsize\color{muted}\llicencia}
\fancyfoot[R]{\small\color{muted}\thepage}
\renewcommand{\headrulewidth}{0.4pt}
\renewcommand{\headrule}{\color{muted!40}\hrule height \headrulewidth}
\fancypagestyle{primera}{\fancyhf{}\renewcommand{\headrulewidth}{0pt}
  \fancyfoot[L]{\scriptsize\color{muted}\llicencia}
  \fancyfoot[R]{\small\color{muted}\thepage}}

% ---------- nivells i exercicis ----------
\newcommand{\estrelles}[1]{{\color{accent}\ifcase#1\or\faStar\faStar[regular]\faStar[regular]\or\faStar\faStar\faStar[regular]\or\faStar\faStar\faStar\fi}}
\newcounter{exer}
\newcommand{\ex}[1]{\par\needspace{4\baselineskip}\medskip\stepcounter{exer}\noindent{\color{navy}\bfseries\theexer.}\ \ {\small\estrelles{#1}}\ \ }

\newcommand{\seccio}[2]{\par\needspace{10\baselineskip}\bigskip
  \begin{tcolorbox}[colback=navy,colframe=navy,coltext=white,arc=2mm,boxrule=0pt,left=3mm,right=3mm,top=1.2mm,bottom=1.2mm]
  {\large\bfseries #1\quad #2}
  \end{tcolorbox}\vspace{-2pt}}

\newcolumntype{C}[1]{>{\centering\arraybackslash}p{#1}}
\usepackage{pifont}

% quadre de recordatori (només a la fitxa d'exercicis)
\newcommand{\recorda}[1]{\par\smallskip\begin{tcolorbox}[colback=accent!7,colframe=accent!55,arc=2mm,boxrule=0.5pt,left=2.5mm,right=2.5mm,top=1mm,bottom=1mm]\small{\color{accent}\faLightbulb[regular]}\ \textbf{Recorda:}\ #1\end{tcolorbox}}
\newcommand{\res}[1]{\boxed{#1}}
\newcommand{\rest}[1]{\fbox{#1}}

\def\fmneed{14}
% =====================================================================
% Capçalera FentMates de les fitxes: un color per a cada curs.
% Els colors són els de la pàgina inicial (LEVELS a assets/js/materials.js);
% el logo ve de _eines/logo/pdf/ (python3 _eines/logo/genera.py).
%
% Ús, al final del preàmbul de la unitat:
%   % =====================================================================
% Capçalera FentMates de les fitxes: un color per a cada curs.
% Els colors són els de la pàgina inicial (LEVELS a assets/js/materials.js);
% el logo ve de _eines/logo/pdf/ (python3 _eines/logo/genera.py).
%
% Ús, al final del preàmbul de la unitat:
%   % =====================================================================
% Capçalera FentMates de les fitxes: un color per a cada curs.
% Els colors són els de la pàgina inicial (LEVELS a assets/js/materials.js);
% el logo ve de _eines/logo/pdf/ (python3 _eines/logo/genera.py).
%
% Ús, al final del preàmbul de la unitat:
%   \input{../comu/fentmates.tex}
%   \fmcurs{1eso}
% i, al començament del document:
%   \capcalera{Matemàtiques}{1r ESO}{Fitxa d'exercicis · Unitat 2}{Divisibilitat}
% Opcional abans de \input: \def\fmneed{14} (espai mínim abans de cada secció).
% =====================================================================
\makeatletter
\def\fm@logodir{../../logo/pdf}
\providecommand{\fmneed}{10}
\def\fm@cursnom{}

% clau del curs, color i color del text damunt del color
\newcommand{\fm@curs}[3]{\@namedef{fm@color@#1}{#2}\@namedef{fm@text@#1}{#3}}
\fm@curs{1eso}{E6B84A}{navy}
\fm@curs{2eso}{3B82F6}{white}
\fm@curs{3eso}{FF7A50}{white}
\fm@curs{4eso-a}{A78BFA}{white}
\fm@curs{4eso-b}{6D28D9}{white}
\fm@curs{1batx-mi}{9B1C1C}{white}
\fm@curs{1batx-ccss}{EF6B6B}{white}
\fm@curs{2batx-mii}{1E6B47}{white}
\fm@curs{2batx-ccss}{2BB3A3}{white}

\newcommand{\fmcurs}[1]{%
  \@ifundefined{fm@color@#1}{\PackageError{fentmates}{Curs desconegut: #1}{}}{}%
  \edef\fm@tmp{\noexpand\definecolor{curs}{HTML}{\@nameuse{fm@color@#1}}%
               \noexpand\colorlet{curstext}{\@nameuse{fm@text@#1}}}\fm@tmp
  \colorlet{cursfosc}{curs!72!black}%
  \def\fm@logo{\fm@logodir/vertical-#1.pdf}%
  \def\fm@marca{\fm@logodir/marca-#1.pdf}}

% ---------- capçalera de la primera pàgina ----------
% #1 matèria (petit), #2 curs, #3 tipus de document i unitat, #4 títol
\newcommand{\capcalera}[4]{%
  \gdef\fm@cursnom{#2}%
  \par\noindent
  \begin{tcolorbox}[colback=curs,colframe=curs,boxrule=0pt,arc=2.2mm,width=3.5cm,height=2.15cm,
      valign=center,left=3mm,right=2mm,top=1mm,bottom=1mm,nobeforeafter,box align=center]
    \color{curstext}{\footnotesize #1}\par\vspace{1pt}{\fontsize{19}{21}\selectfont\bfseries #2\par}%
  \end{tcolorbox}\hspace{5mm}%
  \begin{minipage}[c]{\dimexpr\linewidth-3.5cm-5mm-2.9cm\relax}
    {\color{cursfosc}\small\bfseries\MakeUppercase{#3}}\par\vspace{2pt}
    {\color{navy}\fontsize{20}{23}\selectfont\bfseries #4\par}
  \end{minipage}\hfill
  \begin{minipage}[c]{2.6cm}\raggedleft\includegraphics[width=2.6cm]{\fm@logo}\end{minipage}\par
  \vspace{2pt}}

% ---------- barres de secció del color del curs ----------
\renewcommand{\seccio}[2]{\par\needspace{\fmneed\baselineskip}\bigskip
  \begin{tcolorbox}[colback=curs,colframe=curs,coltext=curstext,arc=2mm,boxrule=0pt,left=3mm,right=3mm,top=1.2mm,bottom=1.2mm]
  {\large\bfseries #1\quad #2}
  \end{tcolorbox}\vspace{-2pt}}

% ---------- capçalera i peu de les pàgines ----------
% primera pàgina: logos institucionals al peu
\fancypagestyle{primera}{\fancyhf{}\renewcommand{\headrulewidth}{0pt}
  \fancyfoot[L]{\raisebox{-0.32\height}{\includegraphics[height=0.78cm]{logo-gva.png}}\hspace{5mm}%
    \raisebox{-0.32\height}{\includegraphics[height=0.78cm]{logo-vermellar.png}}}
  \fancyfoot[R]{\scriptsize\color{muted}\href{https://creativecommons.org/licenses/by-nc-sa/4.0/deed.ca}{\ccbyncsa}\ \ CC BY-NC-SA 4.0 · fentmates.org\qquad\small\thepage}}
% la resta: el símbol a l'esquerra i el curs a la dreta, amb el títol al centre com abans
\fancyhead[L]{\raisebox{-0.3\height}{\includegraphics[height=0.62cm]{\fm@marca}}}
\fancyhead[R]{\small\bfseries\color{cursfosc}\fm@cursnom}
\renewcommand{\headrule}{\color{curs!45}\hrule height \headrulewidth}
\makeatother

%   \fmcurs{1eso}
% i, al començament del document:
%   \capcalera{Matemàtiques}{1r ESO}{Fitxa d'exercicis · Unitat 2}{Divisibilitat}
% Opcional abans de \input: \def\fmneed{14} (espai mínim abans de cada secció).
% =====================================================================
\makeatletter
\def\fm@logodir{../../logo/pdf}
\providecommand{\fmneed}{10}
\def\fm@cursnom{}

% clau del curs, color i color del text damunt del color
\newcommand{\fm@curs}[3]{\@namedef{fm@color@#1}{#2}\@namedef{fm@text@#1}{#3}}
\fm@curs{1eso}{E6B84A}{navy}
\fm@curs{2eso}{3B82F6}{white}
\fm@curs{3eso}{FF7A50}{white}
\fm@curs{4eso-a}{A78BFA}{white}
\fm@curs{4eso-b}{6D28D9}{white}
\fm@curs{1batx-mi}{9B1C1C}{white}
\fm@curs{1batx-ccss}{EF6B6B}{white}
\fm@curs{2batx-mii}{1E6B47}{white}
\fm@curs{2batx-ccss}{2BB3A3}{white}

\newcommand{\fmcurs}[1]{%
  \@ifundefined{fm@color@#1}{\PackageError{fentmates}{Curs desconegut: #1}{}}{}%
  \edef\fm@tmp{\noexpand\definecolor{curs}{HTML}{\@nameuse{fm@color@#1}}%
               \noexpand\colorlet{curstext}{\@nameuse{fm@text@#1}}}\fm@tmp
  \colorlet{cursfosc}{curs!72!black}%
  \def\fm@logo{\fm@logodir/vertical-#1.pdf}%
  \def\fm@marca{\fm@logodir/marca-#1.pdf}}

% ---------- capçalera de la primera pàgina ----------
% #1 matèria (petit), #2 curs, #3 tipus de document i unitat, #4 títol
\newcommand{\capcalera}[4]{%
  \gdef\fm@cursnom{#2}%
  \par\noindent
  \begin{tcolorbox}[colback=curs,colframe=curs,boxrule=0pt,arc=2.2mm,width=3.5cm,height=2.15cm,
      valign=center,left=3mm,right=2mm,top=1mm,bottom=1mm,nobeforeafter,box align=center]
    \color{curstext}{\footnotesize #1}\par\vspace{1pt}{\fontsize{19}{21}\selectfont\bfseries #2\par}%
  \end{tcolorbox}\hspace{5mm}%
  \begin{minipage}[c]{\dimexpr\linewidth-3.5cm-5mm-2.9cm\relax}
    {\color{cursfosc}\small\bfseries\MakeUppercase{#3}}\par\vspace{2pt}
    {\color{navy}\fontsize{20}{23}\selectfont\bfseries #4\par}
  \end{minipage}\hfill
  \begin{minipage}[c]{2.6cm}\raggedleft\includegraphics[width=2.6cm]{\fm@logo}\end{minipage}\par
  \vspace{2pt}}

% ---------- barres de secció del color del curs ----------
\renewcommand{\seccio}[2]{\par\needspace{\fmneed\baselineskip}\bigskip
  \begin{tcolorbox}[colback=curs,colframe=curs,coltext=curstext,arc=2mm,boxrule=0pt,left=3mm,right=3mm,top=1.2mm,bottom=1.2mm]
  {\large\bfseries #1\quad #2}
  \end{tcolorbox}\vspace{-2pt}}

% ---------- capçalera i peu de les pàgines ----------
% primera pàgina: logos institucionals al peu
\fancypagestyle{primera}{\fancyhf{}\renewcommand{\headrulewidth}{0pt}
  \fancyfoot[L]{\raisebox{-0.32\height}{\includegraphics[height=0.78cm]{logo-gva.png}}\hspace{5mm}%
    \raisebox{-0.32\height}{\includegraphics[height=0.78cm]{logo-vermellar.png}}}
  \fancyfoot[R]{\scriptsize\color{muted}\href{https://creativecommons.org/licenses/by-nc-sa/4.0/deed.ca}{\ccbyncsa}\ \ CC BY-NC-SA 4.0 · fentmates.org\qquad\small\thepage}}
% la resta: el símbol a l'esquerra i el curs a la dreta, amb el títol al centre com abans
\fancyhead[L]{\raisebox{-0.3\height}{\includegraphics[height=0.62cm]{\fm@marca}}}
\fancyhead[R]{\small\bfseries\color{cursfosc}\fm@cursnom}
\renewcommand{\headrule}{\color{curs!45}\hrule height \headrulewidth}
\makeatother

%   \fmcurs{1eso}
% i, al començament del document:
%   \capcalera{Matemàtiques}{1r ESO}{Fitxa d'exercicis · Unitat 2}{Divisibilitat}
% Opcional abans de \input: \def\fmneed{14} (espai mínim abans de cada secció).
% =====================================================================
\makeatletter
\def\fm@logodir{../../logo/pdf}
\providecommand{\fmneed}{10}
\def\fm@cursnom{}

% clau del curs, color i color del text damunt del color
\newcommand{\fm@curs}[3]{\@namedef{fm@color@#1}{#2}\@namedef{fm@text@#1}{#3}}
\fm@curs{1eso}{E6B84A}{navy}
\fm@curs{2eso}{3B82F6}{white}
\fm@curs{3eso}{FF7A50}{white}
\fm@curs{4eso-a}{A78BFA}{white}
\fm@curs{4eso-b}{6D28D9}{white}
\fm@curs{1batx-mi}{9B1C1C}{white}
\fm@curs{1batx-ccss}{EF6B6B}{white}
\fm@curs{2batx-mii}{1E6B47}{white}
\fm@curs{2batx-ccss}{2BB3A3}{white}

\newcommand{\fmcurs}[1]{%
  \@ifundefined{fm@color@#1}{\PackageError{fentmates}{Curs desconegut: #1}{}}{}%
  \edef\fm@tmp{\noexpand\definecolor{curs}{HTML}{\@nameuse{fm@color@#1}}%
               \noexpand\colorlet{curstext}{\@nameuse{fm@text@#1}}}\fm@tmp
  \colorlet{cursfosc}{curs!72!black}%
  \def\fm@logo{\fm@logodir/vertical-#1.pdf}%
  \def\fm@marca{\fm@logodir/marca-#1.pdf}}

% ---------- capçalera de la primera pàgina ----------
% #1 matèria (petit), #2 curs, #3 tipus de document i unitat, #4 títol
\newcommand{\capcalera}[4]{%
  \gdef\fm@cursnom{#2}%
  \par\noindent
  \begin{tcolorbox}[colback=curs,colframe=curs,boxrule=0pt,arc=2.2mm,width=3.5cm,height=2.15cm,
      valign=center,left=3mm,right=2mm,top=1mm,bottom=1mm,nobeforeafter,box align=center]
    \color{curstext}{\footnotesize #1}\par\vspace{1pt}{\fontsize{19}{21}\selectfont\bfseries #2\par}%
  \end{tcolorbox}\hspace{5mm}%
  \begin{minipage}[c]{\dimexpr\linewidth-3.5cm-5mm-2.9cm\relax}
    {\color{cursfosc}\small\bfseries\MakeUppercase{#3}}\par\vspace{2pt}
    {\color{navy}\fontsize{20}{23}\selectfont\bfseries #4\par}
  \end{minipage}\hfill
  \begin{minipage}[c]{2.6cm}\raggedleft\includegraphics[width=2.6cm]{\fm@logo}\end{minipage}\par
  \vspace{2pt}}

% ---------- barres de secció del color del curs ----------
\renewcommand{\seccio}[2]{\par\needspace{\fmneed\baselineskip}\bigskip
  \begin{tcolorbox}[colback=curs,colframe=curs,coltext=curstext,arc=2mm,boxrule=0pt,left=3mm,right=3mm,top=1.2mm,bottom=1.2mm]
  {\large\bfseries #1\quad #2}
  \end{tcolorbox}\vspace{-2pt}}

% ---------- capçalera i peu de les pàgines ----------
% primera pàgina: logos institucionals al peu
\fancypagestyle{primera}{\fancyhf{}\renewcommand{\headrulewidth}{0pt}
  \fancyfoot[L]{\raisebox{-0.32\height}{\includegraphics[height=0.78cm]{logo-gva.png}}\hspace{5mm}%
    \raisebox{-0.32\height}{\includegraphics[height=0.78cm]{logo-vermellar.png}}}
  \fancyfoot[R]{\scriptsize\color{muted}\href{https://creativecommons.org/licenses/by-nc-sa/4.0/deed.ca}{\ccbyncsa}\ \ CC BY-NC-SA 4.0 · fentmates.org\qquad\small\thepage}}
% la resta: el símbol a l'esquerra i el curs a la dreta, amb el títol al centre com abans
\fancyhead[L]{\raisebox{-0.3\height}{\includegraphics[height=0.62cm]{\fm@marca}}}
\fancyhead[R]{\small\bfseries\color{cursfosc}\fm@cursnom}
\renewcommand{\headrule}{\color{curs!45}\hrule height \headrulewidth}
\makeatother

\fmcurs{3eso}

\begin{document}
\thispagestyle{primera}

\capcalera{Matemàtiques}{3r ESO}{Solucions de la fitxa · Unitat 10}{Proporcionalitat. Teorema de Tales}
\vspace{8pt}

\begin{tcolorbox}[colback=soft,colframe=soft,arc=2mm,left=3mm,right=3mm,top=1.5mm,bottom=1.5mm]
\small Intenta resoldre cada exercici abans de mirar-ne la solució. Si no t'ix el mateix resultat, repassa els passos o consulta els apunts interactius de la unitat: hi trobaràs una calculadora de proporcions i eines per a practicar el teorema de Tales, la semblança i les escales. Els resultats aproximats estan arrodonits a les centèsimes.
\end{tcolorbox}

% =====================================================
\seccio{1}{Raons i proporcions entre segments}

\ex{1} \begin{tasks}(4)
  \task $\frac{12}{8}=\res{1{,}5}$
  \task $\frac{5}{20}=\res{0{,}25}$
  \task $\frac{60}{45}\approx\res{1{,}33}$
  \task $\frac{20}{25}=\res{0{,}8}$
\end{tasks}

\ex{1} \begin{tasks}(3)
  \task $5x=60\Rightarrow\res{x=12}$
  \task $2x=56\Rightarrow\res{x=28}$
  \task $6x=60\Rightarrow\res{x=10}$
  \task $5x=20\Rightarrow\res{x=4}$
  \task $6x=90\Rightarrow\res{x=15}$
  \task $3x=9\Rightarrow\res{x=3}$
\end{tasks}

\ex{1} \begin{tasks}(2)
  \task $2\cdot6=3\cdot4=12$: \rest{sí}
  \task $3\cdot9=27\neq5\cdot6=30$: \rest{no}
  \task $1{,}5\cdot6=2\cdot4{,}5=9$: \rest{sí}
  \task $4\cdot15=60\neq7\cdot8=56$: \rest{no}
\end{tasks}

\ex{2} a) $\frac46=\frac{10}{x}\Rightarrow x=\frac{6\cdot10}{4}=\res{15\text{ cm}}$. \quad b) $\frac{2{,}5}{3}=\frac{7}{x}\Rightarrow x=\frac{3\cdot7}{2{,}5}=\res{8{,}4\text{ cm}}$.

\ex{2} $\frac35=\frac{12}{x}\Rightarrow3x=60\Rightarrow\res{x=20\text{ cm}}$.

\ex{2} Són $3u$ i $4u$: $7u=56\Rightarrow u=8$. Segments de $\res{24\text{ cm}}$ i $\res{32\text{ cm}}$.

\ex{3} $x^2=a\cdot b$: \ a) $x^2=100\Rightarrow\res{x=10}$; \ b) $x^2=16\Rightarrow\res{x=4}$; \ c) $x^2=18\Rightarrow x=\sqrt{18}\approx\res{4{,}24}$.


% =====================================================
\seccio{2}{El teorema de Tales}

\ex{1} $\frac46=\frac6x\Rightarrow4x=36\Rightarrow\res{x=9}$.

\ex{1} \begin{tasks}(1)
  \task $\frac35=\frac{4{,}8}{x}\Rightarrow\res{\overline{B'C'}=8}$
  \task $\frac64=\frac{x}{5}\Rightarrow\res{\overline{BC}=7{,}5}$
  \task $\frac{4}{10}=\frac{x}{15}\Rightarrow\res{\overline{A'B'}=6}$
\end{tasks}

\ex{2} $\frac x6=\frac48\Rightarrow8x=24\Rightarrow\res{x=3}$.

\ex{2} \begin{tasks}(3)
  \task $\frac23=\frac46$: \rest{sí}
  \task $\frac34=0{,}75\neq\frac45=0{,}8$: \rest{no}
  \task $\frac{1{,}5}{2{,}5}=0{,}6=\frac{2{,}4}{4}$: \rest{sí}
\end{tasks}(Recíproc del teorema de Tales.)

\ex{2} $k=\frac{72}{15+20+25}=\frac{72}{60}=1{,}2$: \ $15\cdot1{,}2=\res{18\text{ m}}$, \ $20\cdot1{,}2=\res{24\text{ m}}$, \ $25\cdot1{,}2=\res{30\text{ m}}$.

\ex{3} $\frac{x}{x+4}=\frac35\Rightarrow5x=3x+12\Rightarrow\res{x=6}$ \ (i $\overline{BC}=10$).

\ex{3} $\frac{\overline{A'B'}}{5}=\frac{\overline{A'C'}}{13}=\frac{19{,}5}{13}=1{,}5$: \ $\res{\overline{A'B'}=7{,}5\text{ cm}}$, \ $\res{\overline{B'C'}=12\text{ cm}}$.


% =====================================================
\seccio{3}{Dividir un segment en parts iguals o proporcionals}

\ex{1} Semirecta des d'$A$, 5 marques iguals, unim la cinquena amb $B$ i tracem paral·leles per les altres quatre. Cada part: $\frac{12}{5}=\res{2{,}4\text{ cm}}$.

\ex{1} \begin{tasks}(1)
  \task $\frac{30}{5}=6$: \ $\res{12\text{ cm i }18\text{ cm}}$
  \task $\frac{24}{6}=4$: \ $\res{4,\ 8\text{ i }12\text{ cm}}$
  \task $\frac{45}{9}=5$: \ $\res{10,\ 15\text{ i }20\text{ cm}}$
\end{tasks}

\ex{2} Semirecta auxiliar, 5 marques iguals, unim la cinquena amb $B$ i tracem la paral·lela per la marca 2. Parts: $\frac{10}{5}=2$, \ $\res{4\text{ cm i }6\text{ cm}}$.

\ex{2} Parts proporcionals a 2 i 5: $\frac{35}{7}=5$; \ $\res{\overline{AP}=10\text{ cm}}$, \ $\res{\overline{PB}=25\text{ cm}}$.

\ex{2} $\frac{3{,}6}{8}=0{,}45$: \ $\res{0{,}45\text{ m}}$, \ $\res{0{,}9\text{ m}}$ i $\res{2{,}25\text{ m}}$.

\ex{3} Parts $3u$ i $7u$: $7u-3u=12\Rightarrow u=3$. \ $\res{9\text{ cm}}$ i $\res{21\text{ cm}}$; \ $\overline{AB}=\res{30\text{ cm}}$.


% =====================================================
\seccio{4}{Triangles en posició de Tales}

\ex{1} $\overline{AB}=4+6=10$; \ $\frac{4}{10}=\frac{x}{15}\Rightarrow\res{x=6}$.

\ex{1} \begin{tasks}(1)
  \task $\frac39=\frac x{12}\Rightarrow\res{\overline{B'C'}=4}$
  \task $\frac58=\frac{7{,}5}{x}\Rightarrow\res{\overline{BC}=12}$
  \task $\frac46=\frac6x\Rightarrow\res{\overline{BC}=9}$
\end{tasks}

\ex{2} $\overline{AB}=10$: \ $\frac{6}{10}=\frac{7{,}2}{x}\Rightarrow\res{x=12}$. \ Tales: $\frac64=\frac9y\Rightarrow\res{y=6}$.

\ex{2} El triangle menut i el gran estan en posició de Tales i les altures també són proporcionals: $\frac{4}{12}=\frac{x}{18}\Rightarrow\res{x=6\text{ cm}}$.

\ex{2} Els triangles $ABC$ i $DEC$ tenen l'angle en $C$ igual (oposats pel vèrtex) i $AB\parallel DE$: \ $\frac{\overline{CE}}{\overline{CA}}=\frac96=1{,}5$. \ $x=8\cdot1{,}5=\res{12}$, \ $y=5\cdot1{,}5=\res{7{,}5}$.

\ex{3} Els triangles $ABC$ i $EDC$ són semblants (angle en $C$ oposat pel vèrtex i angles rectes en $B$ i $D$): \ $\frac{\overline{AB}}{4}=\frac{15}{5}\Rightarrow\res{\overline{AB}=12\text{ m}}$.

\ex{3} $\frac{x}{x+6}=\frac{4}{10}\Rightarrow10x=4x+24\Rightarrow\res{x=4}$.


% =====================================================
\seccio{5}{Figures semblants. Raó de semblança}

\ex{1} \begin{tasks}(2)
  \task $\frac64=\frac96=1{,}5$: \rest{sí}, $k=1{,}5$
  \task $\frac42=2\neq\frac85$: \rest{no}
  \task $\frac63=\frac84=2$: \rest{sí}, $k=2$
  \task \rest{sí} (tots els quadrats ho són), $k=\frac73$
\end{tasks}

\ex{1} a) $\res{10,\ 15\text{ i }17{,}5\text{ cm}}$. \quad b) $\res{1{,}6,\ 2{,}4\text{ i }2{,}8\text{ cm}}$.

\ex{2} $k=\frac{27}{9}=3$; \ costat gran $13\cdot3=\res{39\text{ cm}}$; \ l'àrea es multiplica per $k^2=\res{9}$.

\ex{2} $P'=20\cdot1{,}5=\res{30\text{ cm}}$; \ $A'=27{,}5\cdot1{,}5^2=27{,}5\cdot2{,}25\approx\res{61{,}88\text{ cm}^2}$.

\ex{2} $k^2=\frac{72}{18}=4\Rightarrow\res{k=2}$; \ $P'=16\cdot2=\res{32\text{ cm}}$.

\ex{3} $A=96$; \ $k^2=\frac{600}{96}=6{,}25\Rightarrow k=2{,}5$. \ Dimensions: $\res{20\text{ cm}\times30\text{ cm}}$.

\ex{3} Àrea del menut: $\frac{3\cdot4}{2}=6$; \ $k^2=\frac{54}{6}=9\Rightarrow k=3$. Hipotenusa del menut: 5. \ Costats: $\res{9,\ 12\text{ i }15\text{ cm}}$.


% =====================================================
\seccio{6}{Criteris de semblança de triangles}

\ex{1} \begin{tasks}(2)
  \task \rest{sí}, $k=3$
  \task \rest{no}: $\frac63=\frac84=2\neq\frac{10}{6}$
  \task \rest{sí}, $k=1{,}5$
  \task \rest{no}: $\frac64=\frac96=1{,}5\neq\frac{12}{9}$
\end{tasks}

\ex{1} \begin{tasks}(1)
  \task El primer té 35°, 75°, 70°: \rest{sí} (criteri 1r).
  \task El primer té 50°, 60°, 70° i el segon 60°, 80°, 40°: \rest{no}.
  \task Tots dos tenen 90°, 30° i 60°: \rest{sí}.
\end{tasks}

\ex{2} $\frac64=\frac96=1{,}5$ i l'angle comprés és igual: \rest{sí} (criteri 3r), $k=1{,}5$. \ $\overline{B'C'}=5\cdot1{,}5=\res{7{,}5}$.

\ex{2} $k=\frac96=1{,}5$: \ costats $\res{12\text{ cm i }15\text{ cm}}$; \ raó de les àrees $k^2=\res{2{,}25}$.

\ex{2} Els angles iguals fan $\frac{180-40}{2}=70^\circ$: tots dos tenen 40°, 70°, 70°, i són \rest{semblants}. Dos isòsceles qualssevol, \rest{no} necessàriament.

\ex{2} $a=\sqrt{36+64}=10$. L'altura divideix el triangle en dos triangles semblants al gran: $\frac h8=\frac6{10}\Rightarrow\res{h=4{,}8\text{ cm}}$.

\ex{3} $\overline{BC}=\sqrt{81+144}=15$. Els triangles $HBA$ i $ABC$ són semblants (rectangles amb l'angle $\hat B$ comú): $\frac{\overline{BH}}{9}=\frac{9}{15}\Rightarrow\res{\overline{BH}=5{,}4}$; \ $\res{\overline{HC}=9{,}6}$; \ $\frac{\overline{AH}}{12}=\frac{9}{15}\Rightarrow\res{\overline{AH}=7{,}2}$ cm.

\ex{3} Els triangles $OAB$ i $OCD$ són semblants (angles en $O$ oposats pel vèrtex i $AB\parallel CD$): $\frac{\overline{OC}}{9}=\frac{8}{12}\Rightarrow\res{\overline{OC}=6\text{ cm}}$.


% =====================================================
\seccio{7}{Aplicacions: altures i distàncies}

\ex{1} $\frac{H}{10}=\frac{1}{1{,}25}\Rightarrow\res{H=8\text{ m}}$.

\ex{1} $\frac{s}{1{,}6}=\frac{18}{24}\Rightarrow\res{s=1{,}2\text{ m}}$.

\ex{2} Els dos triangles són semblants: $\frac{H}{18}=\frac{1{,}5}{2}\Rightarrow\res{H=13{,}5\text{ m}}$.

\ex{2} $\frac{H}{146}=\frac22=1\Rightarrow\res{H=146\text{ m}}$ (en aquell moment, l'altura i l'ombra són iguals).

\ex{2} Triangles en posició de Tales: $\frac{H}{30}=\frac{1{,}5}{2}\Rightarrow\res{H=22{,}5\text{ m}}$.

\ex{3} $\frac{1{,}6}{s}=\frac{5}{s+3}\Rightarrow1{,}6s+4{,}8=5s\Rightarrow s=\frac{4{,}8}{3{,}4}\approx\res{1{,}41\text{ m}}$.

\ex{3} Els triangles $ABC$ i $EDC$ són semblants (angles rectes en $B$ i $D$ i angles en $C$ oposats pel vèrtex): $\frac{\overline{AB}}{6}=\frac{20}{8}\Rightarrow\res{\overline{AB}=15\text{ m}}$.

\ex{3} Traçant l'horitzontal dels ulls queden dos triangles en posició de Tales: $\frac{H-1{,}5}{10}=\frac{3-1{,}5}{3}=0{,}5\Rightarrow H-1{,}5=5\Rightarrow\res{H=6{,}5\text{ m}}$.


% =====================================================
\seccio{8}{Escales. Àrees i volums de figures semblants}

\ex{1} \begin{tasks}(3)
  \task $400\text{ cm}=\res{4\text{ m}}$
  \task $150\,000\text{ cm}=\res{1{,}5\text{ km}}$
  \task $700\,000\text{ cm}=\res{7\text{ km}}$
\end{tasks}

\ex{1} a) $1200:100=\res{12\text{ cm}}$. \quad b) $4\,500\,000:1\,500\,000=\res{3\text{ cm}}$.

\ex{2} a) $2\text{ km}=200\,000$ cm; \ $\frac{5}{200\,000}=\frac{1}{40\,000}$: \ $\res{1:40\,000}$. \quad b) $\frac{3}{15\,000}$: \ $\res{1:5000}$.

\ex{2} $3\cdot200=600$ cm $=6$ m; \ $2{,}5\cdot200=500$ cm $=5$ m. \ $\res{6\text{ m}\times5\text{ m}}$, \ $A=\res{30\text{ m}^2}$.

\ex{2} $500\text{ m}=50\,000$ cm; \ $\frac{2}{50\,000}=\frac{1}{25\,000}$: \ $\res{1:25\,000}$. \ $7\cdot25\,000=175\,000$ cm $=\res{1{,}75\text{ km}}$.

\ex{2} $450:18=\res{25\text{ cm}}$. \ $A=40\cdot18^2=40\cdot324=12\,960\text{ cm}^2\approx\res{1{,}30\text{ m}^2}$.

\ex{3} $k=2$: \ capacitat $0{,}25\cdot2^3=\res{2\text{ L}}$; \ etiqueta $60\cdot2^2=\res{240\text{ cm}^2}$.

\ex{3} 1 cm del mapa són $50\,000$ cm $=0{,}5$ km; \ 1 cm² són $0{,}5\cdot0{,}5=0{,}25$ km². \ $6\cdot0{,}25=\res{1{,}5\text{ km}^2}=\res{150\text{ ha}}$.


\end{document}
